% GENERATED FILE. Edit canonical lessons, then run npm run generate:books.
% canonical-source-hash: fnv1a64:82b6c7e3f25c972d
% canonical-lessons: KA-C527-dattu-tegeduko, KA-C527-salahu, KA-C527-posisu, KA-C527-ambegalu, KA-C527-malatayi

\chapter{Raising a Child}
\label{ch:ka-raising-a-child}
\hlchaptermodality{527}

\begin{chapteropening}
\textbf{By the end of this chapter:} \emph{I can say to adopt, to bring up, to nourish, crawling on all fours and a stepmother.}

Complete the last lesson of chapter 527: I can say to adopt, to bring up, to nourish, crawling on all fours and a stepmother.
\end{chapteropening}

\section[\texorpdfstring{\kn{ದತ್ತು} \kn{ತೆಗೆದುಕೊ} (dattu tegeduko)}{ದತ್ತು ತೆಗೆದುಕೊ (dattu tegeduko)}]{\kn{ದತ್ತು} \kn{ತೆಗೆದುಕೊ} (dattu tegeduko) — to adopt}

\label{lesson:KA-C527-dattu-tegeduko}



\begin{quote}
Before the new one: say the Kannada for a great-grandson, then the Kannada for younger.
\end{quote}

\subsection*{You'll want to know: \kn{ದತ್ತು} \kn{ತೆಗೆದುಕೊ}}
\textbf{\kn{ದತ್ತು} \kn{ತೆಗೆದುಕೊ}} — \emph{dattu tegeduko} — \textquotedblleft{}to adopt\textquotedblright{}.

\begin{quote}
\kn{ಆ} \kn{ದಂಪತಿ} \kn{ಒಂದು} \kn{ಹೆಣ್ಣು} \kn{ಮಗುವನ್ನು} \kn{ದತ್ತು} \kn{ತೆಗೆದುಕೊಂಡರು}.
\end{quote}

\begin{quote}
\emph{ā dampati ondu heṇṇu maguvannu dattu tegedukoṇḍaru.}
\end{quote}

\begin{quote}
That couple adopted a baby girl.
\end{quote}

Past (I …): \kn{ದತ್ತು} \kn{ತೆಗೆದುಕೊಂಡೆ}, \emph{dattu tegedukoṇḍe}.

\kn{ದತ್ತು} is adoption, from Sanskrit \kn{ದತ್ತ}, given; \kn{ದತ್ತು} \kn{ಕೊಡು}, to give a child in adoption.

\begin{cousinweb}
\textbf{\kn{ದತ್ತು} \kn{ತೆಗೆದುಕೊ}} (\emph{dattu tegeduko}) joins \emph{dattu}, \textquotedblleft{}adoption\textquotedblright{}, from Sanskrit \emph{datta}, \textquotedblleft{}given\textquotedblright{} (CDIAL 6140), and the native \kn{ತೆಗೆದುಕೊ}, \textquotedblleft{}to take\textquotedblright{}: to take a child given to you.
\end{cousinweb}

\begin{grammarlens}[title={What you've built}]
One more word for describing what you lived and felt.
\end{grammarlens}

\subsection*{Your turn}
\begin{itemize}
\raggedright
  \item \emph{Say it:} \emph{dattu tegeduko}
  \item \emph{Say it:} \emph{dattu tegeduko}, once more
  \item \emph{Say it:} \emph{ā dampati ondu heṇṇu maguvannu dattu tegedukoṇḍaru.} — That couple adopted a baby girl.
  \item \emph{Recall:} say the Kannada for a great-grandson, then the Kannada for younger, then say \emph{dattu tegeduko} again
\end{itemize}

\begin{tcolorbox}[breakable,colback=teal!4,colframe=teal!35!black,title={Before you move on}]
What does \kn{ದತ್ತು} \kn{ತೆಗೆದುಕೊ} mean? (\textquotedblleft{}To adopt\textquotedblright{}.) Say it once more.
\end{tcolorbox}
\section[\texorpdfstring{\kn{ಸಲಹು} (salahu)}{ಸಲಹು (salahu)}]{\kn{ಸಲಹು} (salahu) — to bring up; to protect}

\label{lesson:KA-C527-salahu}



\begin{quote}
Before the new one: say the Kannada for younger, then the Kannada for to adopt.
\end{quote}

\subsection*{You'll want to know: \kn{ಸಲಹು}}
\textbf{\kn{ಸಲಹು}} — \emph{salahu} — \textquotedblleft{}to bring up; to protect\textquotedblright{}.

\begin{quote}
\kn{ಅಜ್ಜಿ} \kn{ತಾಯಿ} \kn{ಇಲ್ಲದ} \kn{ಮಗುವನ್ನು} \kn{ಸಲಹಿದರು}.
\end{quote}

\begin{quote}
\emph{ajji tāyi illada maguvannu salahidaru.}
\end{quote}

\begin{quote}
Grandmother brought up the motherless child.
\end{quote}

Past (I …): \kn{ಸಲಹಿದೆ}, \emph{salahide}.

An old, warm word: to shelter and raise.

\begin{cousinweb}
\textbf{\kn{ಸಲಹು}} (\emph{salahu}) is native Dravidian: DEDR 2781 lists Kannada \emph{salapu}, \emph{salavu}, \emph{salahu}, \textquotedblleft{}to forward, foster, tend, preserve, bring up\textquotedblright{}.
\end{cousinweb}

\begin{grammarlens}[title={What you've built}]
One more word for describing what you lived and felt.
\end{grammarlens}

\subsection*{Your turn}
\begin{itemize}
\raggedright
  \item \emph{Say it:} \emph{salahu}
  \item \emph{Say it:} \emph{salahu}, once more
  \item \emph{Say it:} \emph{ajji tāyi illada maguvannu salahidaru.} — Grandmother brought up the motherless child.
  \item \emph{Recall:} say the Kannada for younger, then the Kannada for to adopt, then say \emph{salahu} again
\end{itemize}

\begin{tcolorbox}[breakable,colback=teal!4,colframe=teal!35!black,title={Before you move on}]
What does \kn{ಸಲಹು} mean? (\textquotedblleft{}To bring up; to protect\textquotedblright{}.) Say it once more.
\end{tcolorbox}
\section[\texorpdfstring{\kn{ಪೋಷಿಸು} (pōṣisu)}{ಪೋಷಿಸು (pōṣisu)}]{\kn{ಪೋಷಿಸು} (pōṣisu) — to nourish; to support}

\label{lesson:KA-C527-posisu}



\begin{quote}
Before the new one: say the Kannada for to adopt, then the Kannada for to bring up.
\end{quote}

\subsection*{You'll want to know: \kn{ಪೋಷಿಸು}}
\textbf{\kn{ಪೋಷಿಸು}} — \emph{pōṣisu} — \textquotedblleft{}to nourish; to support\textquotedblright{}.

\begin{quote}
\kn{ಅವನು} \kn{ತನ್ನ} \kn{ದುಡಿಮೆಯಿಂದ} \kn{ಇಡೀ} \kn{ಕುಟುಂಬವನ್ನು} \kn{ಪೋಷಿಸುತ್ತಾನೆ}.
\end{quote}

\begin{quote}
\emph{avanu tanna duḍimeyinda iḍī kuṭumbavannu pōṣisuttāne.}
\end{quote}

\begin{quote}
He supports the whole family with his earnings.
\end{quote}

Past (I …): \kn{ಪೋಷಿಸಿದೆ}, \emph{pōṣiside}.

\kn{ಪೋಷಕರು}, parents or guardians.

\begin{cousinweb}
\textbf{\kn{ಪೋಷಿಸು}} (\emph{pōṣisu}) is built on Sanskrit \emph{poṣaṇa}, \textquotedblleft{}nourishing, cherishing\textquotedblright{} (Monier-Williams), with the verb ending -\emph{isu}.
\end{cousinweb}

\begin{grammarlens}[title={What you've built}]
One more word for describing what you lived and felt.
\end{grammarlens}

\subsection*{Your turn}
\begin{itemize}
\raggedright
  \item \emph{Say it:} \emph{pōṣisu}
  \item \emph{Say it:} \emph{pōṣisu}, once more
  \item \emph{Say it:} \emph{avanu tanna duḍimeyinda iḍī kuṭumbavannu pōṣisuttāne.} — He supports the whole family with his earnings.
  \item \emph{Recall:} say the Kannada for to adopt, then the Kannada for to bring up, then say \emph{pōṣisu} again
\end{itemize}

\begin{tcolorbox}[breakable,colback=teal!4,colframe=teal!35!black,title={Before you move on}]
What does \kn{ಪೋಷಿಸು} mean? (\textquotedblleft{}To nourish; to support\textquotedblright{}.) Say it once more.
\end{tcolorbox}
\section[\texorpdfstring{\kn{ಅಂಬೆಗಾಲು} (ambegālu)}{ಅಂಬೆಗಾಲು (ambegālu)}]{\kn{ಅಂಬೆಗಾಲು} (ambegālu) — crawling on all fours}

\label{lesson:KA-C527-ambegalu}



\begin{quote}
Before the new one: say the Kannada for to bring up, then the Kannada for to nourish.
\end{quote}

\subsection*{You'll want to know: \kn{ಅಂಬೆಗಾಲು}}
\textbf{\kn{ಅಂಬೆಗಾಲು}} — \emph{ambegālu} — \textquotedblleft{}crawling on all fours\textquotedblright{}.

\begin{quote}
\kn{ಮಗು} \kn{ಅಂಬೆಗಾಲು} \kn{ಇಡುತ್ತಾ} \kn{ಅಡುಗೆಮನೆಗೆ} \kn{ಬಂದಿತು}.
\end{quote}

\begin{quote}
\emph{magu ambegālu iḍuttā aḍugemanege banditu.}
\end{quote}

\begin{quote}
The baby crawled into the kitchen.
\end{quote}

\kn{ಅಂಬೆಗಾಲು} \kn{ಇಡು}, to crawl, of a baby.

\begin{cousinweb}
\textbf{\kn{ಅಂಬೆಗಾಲು}} (\emph{ambegālu}) is native Dravidian: DEDR 180 lists Kannada \emph{ambegālu}, \textquotedblleft{}the hands and knees placed on the ground for crawling\textquotedblright{}.
\end{cousinweb}

\begin{grammarlens}[title={What you've built}]
One more word for describing what you lived and felt.
\end{grammarlens}

\subsection*{Your turn}
\begin{itemize}
\raggedright
  \item \emph{Say it:} \emph{ambegālu}
  \item \emph{Say it:} \emph{ambegālu}, once more
  \item \emph{Say it:} \emph{magu ambegālu iḍuttā aḍugemanege banditu.} — The baby crawled into the kitchen.
  \item \emph{Recall:} say the Kannada for to bring up, then the Kannada for to nourish, then say \emph{ambegālu} again
\end{itemize}

\begin{tcolorbox}[breakable,colback=teal!4,colframe=teal!35!black,title={Before you move on}]
What does \kn{ಅಂಬೆಗಾಲು} mean? (\textquotedblleft{}Crawling on all fours\textquotedblright{}.) Say it once more.
\end{tcolorbox}
\section[\texorpdfstring{\kn{ಮಲತಾಯಿ} (malatāyi)}{ಮಲತಾಯಿ (malatāyi)}]{\kn{ಮಲತಾಯಿ} (malatāyi) — a stepmother}

\label{lesson:KA-C527-malatayi}



\begin{quote}
Before the new one: say the Kannada for to nourish, then the Kannada for crawling on all fours.
\end{quote}

\subsection*{You'll want to know: \kn{ಮಲತಾಯಿ}}
\textbf{\kn{ಮಲತಾಯಿ}} — \emph{malatāyi} — \textquotedblleft{}a stepmother\textquotedblright{}.

\begin{quote}
\kn{ಆ} \kn{ಕಥೆಯಲ್ಲಿ} \kn{ಮಲತಾಯಿ} \kn{ತುಂಬಾ} \kn{ಒಳ್ಳೆಯವಳು}.
\end{quote}

\begin{quote}
\emph{ā katheyalli malatāyi tumbā oḷḷeyavaḷu.}
\end{quote}

\begin{quote}
In that story the stepmother is very kind.
\end{quote}

\kn{ಮಲ}- marks a step-relation: \kn{ಮಲತಮ್ಮ}, a younger half-brother.

\begin{cousinweb}
\textbf{\kn{ಮಲತಾಯಿ}} (\emph{malatāyi}) is a Kannada compound: \emph{mala-}, "step-", whose origin is not settled, and \emph{tāyi}, \textquotedblleft{}mother\textquotedblright{} (DEDR 364).
\end{cousinweb}

\begin{grammarlens}[title={What you've built}]
One more word for describing what you lived and felt.
\end{grammarlens}

\subsection*{Your turn}
\begin{itemize}
\raggedright
  \item \emph{Say it:} \emph{malatāyi}
  \item \emph{Say it:} \emph{malatāyi}, once more
  \item \emph{Say it:} \emph{ā katheyalli malatāyi tumbā oḷḷeyavaḷu.} — In that story the stepmother is very kind.
  \item \emph{Recall:} say the Kannada for to nourish, then the Kannada for crawling on all fours, then say \emph{malatāyi} again
\end{itemize}

\begin{tcolorbox}[breakable,colback=teal!4,colframe=teal!35!black,title={Before you move on}]
What does \kn{ಮಲತಾಯಿ} mean? (\textquotedblleft{}A stepmother\textquotedblright{}.) Say it once more.
\end{tcolorbox}
